\chapter{Conclusion}
\label{sec:conclusion}

\section{Summary of the Example}
This demonstration has described an illustrative comparison between a baseline
retrieval configuration and a reranked variant. The proposed procedure holds
the collection and queries fixed, retains paired query scores and examines
both overall and group summaries. It is a worked example of thesis structure,
not evidence that a particular retrieval method outperforms another.

The document connects research questions to an experimental plan and then to
a discussion of possible results. That connection is useful in a completed
thesis because it lets the reader assess whether the reported evidence answers
the questions posed at the beginning. A conclusion should revisit those
questions directly instead of introducing an unrelated set of claims.

\section{Reporting Contributions}
A final thesis should distinguish a software artifact, a methodological
contribution and an empirical finding. Implementing a system establishes that
an artifact exists; evaluating it under a documented protocol provides a
different kind of evidence. Neither automatically demonstrates that the
approach generalises beyond the setting that was studied.

For this template, the concrete artifact is the formatted document itself.
It demonstrates chapter openings with visible page numbers, continuation
pages with running headers, and uninterrupted chapter transitions on either
page parity. Tables, a figure, citations and an appendix provide additional
examples for reviewing the surrounding document structure.

% Demonstration-only break.
\newpage
\section{Boundaries of the Conclusions}
A conclusion should state the conditions under which its claims hold. Those
conditions may include the collection version, query population, relevance
judgements and computational constraints. Reporting them does not weaken an
otherwise sound result; it helps the reader understand what would need to
remain comparable when applying the finding elsewhere.

The invented scores in this document support no empirical claim. Their role
is to exercise table alignment, decimal formatting and cross-references.
Similarly, the fictional references demonstrate bibliography syntax without
establishing a literature basis for the study. Both should be replaced before
the template is used to report actual research.

\section{Closing Remarks}
The final paragraph of a thesis can explain why the established contribution
matters while keeping its scope precise. It should follow from the evidence
already presented rather than rely on a promise that future work will resolve
all remaining issues. Specific possible extensions belong in the next chapter,
where their motivation and requirements can be explained separately.

For layout inspection, compare this page with the preceding chapter-opening
page. Both should display the same centred footer style. Only this continuation
page should have a running header and a thin horizontal rule above the text.
The next chapter starts on the next page, without an intervening empty leaf.
